\chapter{Design}
\label{chap:design}

In the previous chapter, we covered two BDD libraries that utilize parallelism in the \texttt{apply} operation (described in detail in Section \ref{chap:bdd:apply}). We focused mainly on the OxiDD framework. This is mainly due
to its implementation in the Rust programming language, its inspiration from Sylvan, and how relatively new it is. In
this chapter, we will go through the design of a BDD library that uses breadth-first-search parallelism. However, the
scope of this library is focused on \texttt{apply} and its scalability. Therefore, the product should not be a
full-fledged user-facing library, but should satisfy the following requirements:

\begin{itemize}
    \item The \texttt{apply} operation distributes the work among multiple worker threads in a breadth-first-search
        manner.
    \item The implementation is thoroughly tested to guarantee its correctness.
    \item The resulting library is benchmarked on BDDs of diverse sizes and across various thread counts (including
        single-threaded) to assess the usefulness and scalability of the BFS approach.
\end{itemize}

The features commonly implemented in other BDD libraries, and specifically out of scope, include evaluating the BDD
for a given valuation of variables, providing a user API to create variables and new functions, and checking for
satisfiability. The design discussed in this chapter should be seen as a foundation that can be used in the future
either as a learning resource or as a baseline for further optimizations.


\section{Historical context}
\label{chap:design:history}
Several times in this text, a paper \cite{bfs-bdd} by Ochi, H., Yasuoka, K., and Yajima, S. was cited. The name would
suggest it could provide a solid foundation for this work. However, it was published in 1993, so the problem it
solved was memory capacity. Large BDDs had difficulty fitting in main memory (commonly referred to as RAM). The paper
specifically mentions the Sun SPARCstation 10 workstation, which had 64 MB of RAM. The researchers proposed a 1 GB
virtual space that combined RAM with secondary storage on a 1 GB hard disk drive. Their solution lies in organizing
all nodes and hash tables strictly by variable level in contiguous blocks on disk. The nodes in the same level are
loaded from disk to main memory, processed, and written back to disk \cite{bfs-bdd}.

In today's age, even low-end personal laptops have at least 4 GB of RAM, let alone professional workstations used for
BDD computations, which are equipped with hundreds of GB of main memory. Although memory consumption concerns are not
frequently addressed at this time, the principles of breadth-first execution have not changed.


\section{General ideas}
\label{chap:design:general}

The very first design decision was to choose the type of BDD representation --- split or shared\footnote{The
definition is covered in Section \ref{chap:other-impls:oxidd}.}. Since memory optimizations are taken into
consideration, but not the main focus of the work, the shared representation is the more reasonable choice for its
memory savings, and its implementation is relatively easy.

The cornerstone of all data structures is the representation of BDD primitives, such as a Node. Listing
\ref{lst:key-definitions} depicts structures \texttt{NodeId}, \texttt{Variable}, and \texttt{Node}. For simplicity's
sake, the visibility modifiers and attribute macros are omitted. This representation is inspired by the LibBDD
\footnote{\url{https://github.com/sybila/biodivine-lib-bdd}} library.

\begin{listing}[htbp]
    \caption{Definitions of the key types --- \texttt{NodeId}, \texttt{Variable}, and \texttt{Node}.}
    \label{lst:key-definitions}
    \begin{minted}{rust}
    struct NodeId(Uuid);
    struct Variable(u32);
    struct Node {
        variable: Variable,
        low_child: NodeId,
        high_child: NodeId,
    }
    \end{minted}
\end{listing}

Each node has an assigned variable and references to its successors (hereafter referred to as children). Rather than
direct references, these identifiers act as keys into the library's node hash map. You may notice that
\texttt{NodeId} and \texttt{Variable} use the \texttt{struct} keyword instead of \texttt{type}. This pattern,
called Newtype, is useful because the new type can expose a public API that is different from the API of the
private inner type \cite[ch. 20]{rust-book}.

- representation of the BDD

- BFS - we need a queue and hashtable, justify using crossbeam-queue and DashMap

- General overview was covered in \ref{chap:bdd:apply:iterative}

- for apply, just a single operation is needed - logical and. Other operations should be farely trivial to implement

- The termination detection for generating tasks

- Identifying tasks after deduplication

- ensure node - two maps that are basically identical, but key,value are swapped

